\documentclass[11pt]{article}

\usepackage[utf8]{inputenc}
\usepackage[T1]{fontenc}
\usepackage{amsmath,amssymb,amsthm}
\usepackage{booktabs}
\usepackage{hyperref}
\usepackage{natbib}
\usepackage[margin=1in]{geometry}

\newcommand{\credence}{\textsc{Credence}}

\title{Cheap Design for LLM Agents: \\An Information-Geometric Perspective}

\author{
  Guy Freeman \\
  Independent Researcher, Tel Aviv
}

\date{\today}

\begin{document}
\maketitle

\begin{abstract}
% Core claim: The Credence decision layer is an embodied universal approximator in the
% sense of Ay [2015] for the constraint set defined by tool reliability profiles,
% achieving the full behavioural repertoire with exponentially fewer parameters than a
% universal approximator (the LLM).
\end{abstract}


% ============================================================
% EXTRACTED FROM PAPER 1:
% ============================================================

% --- From §2.3 "Embodied agent design" (Paper 1 line 97) ---
%
% \paragraph{Embodied agent design.} \citet{ay2015geometric} formalises the sensorimotor loop as a causal model with Markov kernels and shows that an agent's controller need not be a universal approximator---it suffices to match the behavioural repertoire afforded by its embodiment. This \emph{cheap design} principle \citep{pfeifer2006body, montufar2015cheap} motivates \credence{}'s architecture: the Bayesian decision layer is far simpler than the LLM it coordinates, yet sufficient for optimal tool selection given the embodiment constraints (available tools and their reliability profiles).

% --- From §3.2 Architecture: Brain, Body, Environment (Paper 1 lines 137-148) ---
%
% \subsection{Architecture: Brain, Body, Environment}
% \label{sec:architecture}
%
% Following \citet{ay2015geometric}, we decompose the agent into three components:
%
% \begin{itemize}
% \item \textbf{Brain} (controller $\mathcal{C}$): The Bayesian decision layer. Maintains beliefs about tool reliability, computes expected utilities, selects actions. Implemented in ${\sim}200$ lines of Python/NumPy. Contains no LLM.
% \item \textbf{Body} (actuators and sensors): The available tools, including the LLM itself. These are \emph{noisy channels} in the information-theoretic sense---Markov kernels $\beta: \mathcal{W} \to \Delta_\mathcal{S}$ from world states to sensor states, following the notation of \citet{ay2015geometric}. The LLM is a prosthetic: a powerful but unreliable sensor whose reliability profile is learned, not assumed.
% \item \textbf{Environment} (world $\mathcal{W}$): The questions, their categories, and the ground-truth answers. The agent does not observe the environment directly---only through the noisy responses of its tools.
% \end{itemize}
%
% This separation instantiates the principle of cheap design: the brain is deliberately simple because the body (LLM + tools) provides the perceptual complexity. The brain's job is not to understand language or reason about content---it is to decide \emph{which tool to query and when to stop}.

% --- From §6 "Cheap design for LLM agents" (Paper 1 lines 382-387) ---
%
% \paragraph{Cheap design for LLM agents.}
% \citet{ay2015geometric} argues that an embodied agent's controller need not be a universal approximator---it need only be sufficient for expressing desired behaviours given the agent's embodiment constraints. The \credence{} decision layer demonstrates this principle: a Julia inference engine maintaining Beta distributions and computing VOI outperforms a full LLM-based decision-maker that is orders of magnitude more complex. The LLM is powerful but expensive; using it for decisions it is provably ill-suited to make---resource allocation under uncertainty---is the opposite of cheap design.
%
% The analogy to Ay's exponential family policy models (\citealt{ay2015geometric}, Proposition~1) is suggestive: just as Ay shows that policies parameterised by a small number of ``feature vectors'' derived from the embodiment constraints can be embodied universal approximators, the \credence{} decision layer achieves near-oracle performance using only the sufficient statistics of the Beta posteriors. The ``feature vectors'' are the tool-category reliability estimates; the ``embodiment constraints'' are the tools' true (unknown) reliability profiles.
%
% The computational asymmetry is stark: the \credence{} decision layer processes each question in $0.07$ seconds (closed-form Beta posterior updates and VOI computation), whilst the best LLM agent requires $5.65$ seconds---an $80\times$ difference driven entirely by LLM inference overhead for tool selection decisions that admit analytical solutions.


% ============================================================
% NEW MATERIAL NEEDED:
% ============================================================

% 1. Explicit kernel correspondence: map Credence components onto Ay's
%    (beta, pi, alpha, phi) Markov kernels
%
% 2. Define the constraint set Sigma for tool reliability profiles
%
% 3. State and prove: "The Beta-posterior VOI controller is a sufficient model
%    (Ay Definition 1) for Sigma"
%
% 4. Connect to Ay's Proposition 1: the exponential family parameterisation maps
%    onto softmax tool selection; the feature vectors are tool-category reliability
%    profiles; dimensionality reduction from |C|*|A| to T*C*2
%
% 5. Connect minimum-entropy result (Ay eq. 17) to empirical finding that Credence
%    uses fewer tools
%
% 6. Discuss controlled RBM (Ay Proposition 4) as suggestive for the full Credence
%    DSL architecture: expression trees as hidden nodes, CRP clustering as learning
%    interaction structure
%
% 7. Address limits: Ay assumes finite state sets; Credence has continuous belief
%    states. Ay doesn't address learning Sigma; Credence learns it online.
%
% 8. Stretch goal: Contact Ay, Montufar, or Zahedi about collaboration.


% ============================================================
% BIBLIOGRAPHY
% ============================================================
\bibliographystyle{plainnat}

\begin{thebibliography}{99}

\bibitem[Ay(2015)]{ay2015geometric}
N.~Ay.
\newblock Geometric design principles for brains of embodied agents.
\newblock \emph{K{\"u}nstliche Intelligenz}, 29(4):389--399, 2015.

\bibitem[Mont\'{u}far et~al.(2015)]{montufar2015cheap}
G.~Mont\'{u}far, K.~Ghazi-Zahedi, and N.~Ay.
\newblock A theory of cheap control in embodied systems.
\newblock \emph{PLOS Computational Biology}, 11(9):e1004427, 2015.

\bibitem[Pfeifer and Bongard(2006)]{pfeifer2006body}
R.~Pfeifer and J.~C. Bongard.
\newblock \emph{How the Body Shapes the Way We Think}.
\newblock MIT Press, 2006.

\end{thebibliography}

\end{document}
